\chapter{Floating-Point Unit (fpu)}
\label{ch:fpu}

\section{Introduction}

\sig{fpu} computes the addition, subtraction, multiplication or division of two
operands, selected by \sig{funct}. It holds the \sig{fadd\_core}, \sig{fmul\_core} and \sig{fdiv\_core}
datapaths of Chapters~\ref{ch:fadd}, \ref{ch:fmul} and~\ref{ch:fdiv}, and they
share the \sig{unpack} and \sig{pack} modules, \sig{pack} with its \sig{norm} stage, as
Figure~\ref{fig:fpu} shows. Subtraction uses \sig{fadd\_core} with the sign of
\sig{op\_b} inverted.

\begin{figure}[H]
\centering
\begin{tikzpicture}[node distance=5mm and 9mm]
  \node[ugblock] (u) {2 $\times$ \sig{unpack}};
  \node[ugblock,right=of u] (m) {\sig{fmul\_core}};
  \node[ugblock,above=of m] (a) {\sig{fadd\_core}};
  \node[ugblock,below=of m] (d) {\sig{fdiv\_core}};
  \node[ugblock,right=of m] (p) {\sig{pack}, with \sig{norm}};
  \draw[ugarrow] ($(u.west)-(9mm,0)$) node[above right,font=\footnotesize\ttfamily,xshift=-1mm] {op} -- (u);
  \draw[ugarrow] (u.east) -- (m.west);
  \draw[ugarrow] (u.east) -- ++(3mm,0) |- (a.west);
  \draw[ugarrow] (u.east) -- ++(3mm,0) |- (d.west);
  \draw[ugarrow] (m.east) -- (p.west);
  \draw[ugarrow] (a.east) -- ++(3mm,0) |- (p.west);
  \draw[ugarrow] (d.east) -- ++(3mm,0) |- (p.west);
  \draw[ugarrow] (p) -- ($(p.east)+(9mm,0)$) node[above left,font=\footnotesize\ttfamily,xshift=1mm] {res};
\end{tikzpicture}
\caption{\sig{fpu} Block Diagram}
\label{fig:fpu}
\end{figure}

\section{Features}

\begin{itemize}
\item Addition, subtraction, multiplication and division in one unit.
\item The latencies of Table~\ref{tab:fpu_funct}.
\item One new operation per clock cycle for any mix of additions, subtractions
  and multiplications, with results in issue order.
\item Subnormal operands and results.
\item The five rounding modes of Table~\ref{tab:rnd_mode}.
\item NaN, infinity and zero-divisor handling, selected by \sig{SPECIAL\_CASES}.
\end{itemize}

\section{Parameters}

\sig{fpu} has the parameters of Table~\ref{tab:common_params}.

\section{Interface Signals}

\sig{fpu} has the signals of Table~\ref{tab:common_sigs} and the signal of
Table~\ref{tab:fpu_sigs}.

\begin{sigtable}{\sig{fpu} Additional Signal}{tab:fpu_sigs}
funct & Input & 2 & Operation, encoded as in Table~\ref{tab:fpu_funct}. \\
\end{sigtable}

\Needspace*{23\baselineskip}
\section{Instantiation}

The example below instantiates \sig{fpu} for binary32. Compile it with the files
of Table~\ref{tab:files}.

\begin{vcode}
fpu #(
   .DATA_W       (32),
   .EXP_W        (8),
   .SPECIAL_CASES(1)
) fpu0 (
   .clk     (clk),
   .rst     (rst),
   .valid_i (valid_i),
   .valid_o (valid_o),
   .rnd_mode(rnd_mode),
   .funct   (funct),
   .op_a    (op_a),
   .op_b    (op_b),
   .res     (res)
);
\end{vcode}

\section{Operation}

\sig{fpu} follows the protocol of Section~\ref{sec:protocol}, with the latencies
of Table~\ref{tab:fpu_funct}. They were measured for the 16-, 32- and 64-bit
formats.

\begin{threetable}{\sig{fpu} Operations and Latencies}{tab:fpu_funct}{\sig{funct}}{Operation}{Latency in clock cycles}{1.8cm}{3.4cm}
\sig{00} & Addition & 5 \\ \hline
\sig{01} & Subtraction, $\mathtt{op\_a}-\mathtt{op\_b}$ & 5 \\ \hline
\sig{10} & Multiplication & 5 \\ \hline
\sig{11} & Division, $\mathtt{op\_a}/\mathtt{op\_b}$ & $\mathtt{MAN\_W}+6$ (30 for binary32) \\
\end{threetable}

\begin{itemize}
\item \textbf{Function select.} \sig{funct} is used in the cycle with
  \sig{valid\_i} high and can change after it.
\item \textbf{Addition, subtraction and multiplication.} These operations have
  equal latency, because \sig{fpu} sets \sig{OUT\_REG} to 1 in \sig{fmul\_core}. They can
  be mixed in any order on consecutive clock cycles. The unit samples their
  operands in the \sig{valid\_i} cycle, and the results return in issue order.
\item \textbf{Division.} A division can start right after any operation, and
  its result returns after the results of the earlier operations. The unit samples
  its operands in the \sig{valid\_i} cycle. Issue nothing else during the
  division.
\item \textbf{After a division.} A new operation, with any \sig{funct}, can start
  in the cycle in which \sig{valid\_o} of the division is high.
\end{itemize}

For special operands, each operation gives the results of its unit:
Table~\ref{tab:fadd_special} for addition and subtraction (with the sign of
\sig{op\_b} inverted for subtraction), Table~\ref{tab:fmul_special} for
multiplication and Table~\ref{tab:fdiv_special} for division.
